\section{Finite-horizon planning with admissible actions}
\label{sec:planning}

An artifact decision changes both what is available and what can be afforded.
Inspection can turn an untried cached artifact into a ready artifact, but a
rejection also spends its inspection fee. Generation can create a candidate,
but deployment must wait for a checked success. A controller therefore needs
an action set that depends on the current state, rather than a fixed menu with
inadmissible actions assigned convenient scores.

\begin{definition}[Guarded rational system]
Let $S$ be a state space, $A(s)$ a finite nonempty set of actions, and $O$ a
finite observation space. A guarded system consists of a rational reward
$r(s,a)$, a transition $T(s,a,o)$, and rational probabilities $P(o\mid s,a)$
satisfying
\[
 P(o\mid s,a)\geq0,\qquad \sum_{o\in O}P(o\mid s,a)=1.
\]
Its finite-horizon Bellman values are
\begin{align}
 V_0(s)&=0,\label{eq:bellmanzero}\\
 Q_h(s,a)&=r(s,a)+\sum_oP(o\mid s,a)V_h(T(s,a,o)),\label{eq:qvalue}\\
 V_{h+1}(s)&=\max_{a\in A(s)}Q_h(s,a).\label{eq:bellman}
\end{align}
\end{definition}

Here $h$ in $Q_h$ counts \emph{continuation} steps. Choosing an acquisition
action and then deploying a successful artifact requires two total actions,
so its first-action score is $Q_1$, not $Q_0$. This indexing also applies to
\texttt{robustChoice md h s} in the formal source.

The state space need not be finite. Rational Bayesian beliefs are one intended
instance. Only each feasible action set and the observation space are finite.
The finite maximum exists even when there is no asserted efficient procedure
for finding it. The mathematical selector and the executable evaluator are
separate constructions.

\begin{theorem}[Optimality over complete adaptive plans]
Every horizon $H$ has a legal deterministic observation tree attaining
$V_H(s)$. Every other legal deterministic observation tree has value at most
$V_H(s)$. If each admissible transition preserves an invariant, the attaining
tree preserves that invariant on every branch.
\end{theorem}
\begin{proof}
At horizon zero, the empty tree has value zero. At horizon $h+1$, choose an
action attaining the finite maximum in \eqref{eq:bellman}; after each observation,
attach an optimal horizon-$h$ tree at the successor state. Induction proves
equality of the resulting expected value and the Bellman value. For an arbitrary
legal tree, induction bounds each child value by the corresponding optimal
continuation. Multiplication by nonnegative probabilities, summation, and the
maximum over admissible first actions give the upper bound. Invariant
preservation follows by the same tree induction. The invariant argument applies
to every observation branch, including a branch of probability zero.
\end{proof}

This is a prose reconstruction of
\path{safe_optimal_plan_exists},
\path{optimalPlan_value}, and \path{planValue_le_value} in
\path{AVOC/GuardedControl.lean}. The competitors are complete adaptive
observation trees. They need not use the same subsequent choices as the selected
controller. The comparator nevertheless uses the same state, information,
action guards, and transition model. An oracle with additional information is a
different optimization problem.

\subsection{Warm inspection and cold generation}

The lazy staged controller makes this structure concrete. Its state records
warm and cold statuses, remaining budget, reserved spending, original credit,
and a terminal flag. Each artifact status is untried, failed, or ready. The
action list always contains stop. Inspection and generation require an untried
status and enough budget. Deployment requires the matching ready status.
Stopping or deploying ends the episode.

An acquisition transition pays its setup fee whether its Boolean report is true
or false. A true report changes the corresponding status to ready. A false
report changes it to failed. Acquisition rewards are negative setup fees,
deployment rewards are nonnegative gross values, and stopping has reward zero.
This prevents three easy accounting mistakes: assuming unsuccessful acquisition
is free, treating a quote as an already available artifact, and deploying before
checked success.

\begin{example}[The artifact-backed rank-one comparison]
The exact quote model retained in \path{AVOC/RobustQuoteExamples.lean} has
warm setup $133$, cold setup $196$, success probability $1$ for both, and gross
deployment value $15000$ for each over $1000$ calls. With initial credit $196$
and two total actions available,
\[
 Q_1(\mathrm{inspectWarm})=15000-133=14867,
 \qquad
 Q_1(\mathrm{generateCold})=15000-196=14804.
\]
Warm inspection is better by $63$ in this exact model. At one total action,
acquisition has only its negative immediate reward, and stopping is preferable
when neither artifact is ready. The deployment opportunity creates the change
between these two horizons.
\end{example}

This example has actual bounded artifact witnesses in the retained source.
The gross value is the symbolic execution saving for the protected workload,
not a measured dollar return or a wall-clock benchmark. The singleton evidence
makes all quote uncertainty radii vanish, so the robust lower-Q selector
recovers the inherited exact selector. The relevant declarations are
\path{exact_recovers_v037},
\path{exact_all_local_radii_zero}, and
\path{artifact_refines_deterministic_fields}.

\subsection{An executable evaluator and its cost boundary}

The executable planner enumerates admissible first actions and observation
branches, recursively evaluates continuations, and chooses a best score. Under
the formal \path{Realizes} relation, its returned score equals the mathematical
Bellman value. A full catalogue enumeration and a verified admissibility check
connect the staged workflow to that evaluator. In particular,
\path{WorkflowPlanner.model_realizes} and
\path{executable_policy_source_safe} connect actual staged transitions,
source-bound artifacts, and the selected policy.

The planner's work counter counts defined recursive evaluator nodes. A branching
cap permits reserving an abstract planning allowance, and truncated lookahead
has a source-backed policy regret bound. This counter is not a rational-bit
complexity bound. It does not include all host parsing, catalogue scans,
reservation arithmetic, storage operations, or physical invocation cost.
Consequently, a proved allowance in that model cannot by itself establish a
time bound for the complete service.

\section{Quote uncertainty and certified lower-Q selection}
\label{sec:quotes}

Quotes carry more than a success estimate. A protected identity records the
artifact kind, exact source signature, workload call count, cost model, provider,
checker, environment, and schema. A setup band has natural endpoints. Gross
and probability bands have rational endpoints, with probability endpoints in
$[0,1]$ and nonnegative gross lower endpoints. A submitted quote validates only
when it equals the protected approved quote. Changing either identity or any
band changes the evidence.

For kind $k\in\{w,c\}$, write
\[
 c_k\in[c_k^-,c_k^+],\qquad
 g_k\in[g_k^-,g_k^+],\qquad
 p_k\in[p_k^-,p_k^+].
\]
The formal nominal model reserves $c_k^+$, credits $g_k^-$, and uses $p_k^-$.
A compatible realization supplies exact values inside these bands. In the true
comparison system, rewards and report probabilities use those exact values,
but the state transitions and action guards remain those of the nominal model.
In particular, setup reservations continue to use the upper setup bound.

This common guard matters. If the true setup is $1$ and the upper quoted setup
is $3$, a budget of $2$ still makes the acquisition inadmissible for both systems
in the theorem. The comparator does not obtain an extra feasible action by
knowing the smaller true fee. A theorem against that larger feasible set would
require a different comparison argument.

\subsection{From primitive bands to value bounds}

Let $w_c(k)=c_k^+-c_k^-$, $w_g(k)=g_k^+-g_k^-$, and
$w_p(k)=p_k^+-p_k^-$. The source uses the uniform quantities
\begin{align}
 \varepsilon&=w_c(w)+w_c(c)+w_g(w)+w_g(c),\\
 \delta&=2\bigl(w_p(w)+w_p(c)\bigr),\\
 M&=c_w^++c_c^++g_w^++g_c^+.
\end{align}
They bound reward error, Boolean-row $L^1$ error, and nominal one-step reward
magnitude respectively. For an acquisition of kind $k$, the sharper local
reward error is $w_c(k)$ and its row error is $2w_p(k)$. For deployment of kind
$k$, the local reward error is $w_g(k)$ and its row error is zero. Stop has zero
primitive errors.

\begin{lemma}[Bernoulli row error and expectation transport]
For Bernoulli parameters $p,q$,
\[
 \sum_{o\in\{0,1\}}|P_p(o)-P_q(o)|=2|p-q|.
\]
If $|W(t)|\leq B$, then replacing an observation row $P$ by $\widehat P$
changes its expected continuation by at most
\[
 \left|\sum_o(P(o)-\widehat P(o))W(T(s,a,o))\right|
 \leq B\sum_o|P(o)-\widehat P(o)|.
\]
\end{lemma}
\begin{proof}
The two Bernoulli row differences are $p-q$ and $q-p$, which have equal
absolute values. For the second statement, apply the triangle inequality to the
finite sum and use the uniform bound on each continuation value.
\end{proof}

Induction on the Bellman recursion gives $|\widehat V_h(s)|\leq hM$ for the
nominal value. The source's accumulated value radius and local Q radius are
\begin{align}
 D_h&=\sum_{j=0}^{h-1}(\varepsilon+\delta jM),\label{eq:value-radius}\\
 \rho_h(a)&=D_h+\varepsilon_a+\delta_a hM,\label{eq:local-radius}\\
 \Gamma_h&=D_h+\varepsilon+\delta hM.\label{eq:global-radius}
\end{align}
The Bellman residual is zero here because the nominal tables are exact Bellman
values. The broader model-error framework also allows a supplied table with a
proved nonzero residual, but no residual term is silently introduced in these
quote results.

\begin{theorem}[Value and Q enclosures]
For every compatible realization, every state, and every continuation horizon,
\begin{align}
 |V_h(s)-\widehat V_h(s)|&\leq D_h,\\
 |Q_h(s,a)-\widehat Q_h(s,a)|&\leq\rho_h(a)\leq\Gamma_h.
\end{align}
Thus $L_h(s,a)=\widehat Q_h(s,a)-\rho_h(a)$ and
$U_h(s,a)=\widehat Q_h(s,a)+\rho_h(a)$ enclose the true Q value.
\end{theorem}
\begin{proof}
At horizon zero the value error is zero. In a single backup, insert and subtract
the expectation of the nominal continuation under the true row. The reward
error contributes at most $\varepsilon_a$, the continuation-value error at most
$D_h$, and row transport at most $\delta_a hM$. Taking a maximum over the same
feasible action set is nonexpansive in the uniform score error. Therefore
$D_{h+1}\leq D_h+\varepsilon+\delta hM$, whose accumulated upper bound is
\eqref{eq:value-radius}. The local Q statement follows from the same inserted
expectation before taking a maximum.
\end{proof}

The reconstructed argument follows \path{value_error},
\path{local_q_error}, \path{value_between}, and
\path{q_between} in \path{AVOC/RobustQuotePlanning.lean}. The primitive
steps are \path{reward_close}, \path{probability_l1_le}, and
\path{expectation_transport}.

\subsection{What the robust selector optimizes}

The executable robust choice maximizes $L_h(s,a)$ over feasible actions. It is
a maximizer of a \emph{certified surrogate lower bound}. It is not asserted to
solve the exact minimax problem over all compatible realizations. A bound can
be valid and substantially below the actual worst-case score.

\begin{proposition}[Distinct-competitor dominance]
If $a$ is feasible and
\[
 U_h(s,b)\leq L_h(s,a)\quad
 \text{for every feasible }b\ne a,
\]
then $a$ is optimal in every compatible realization.
\end{proposition}
\begin{proof}
For a distinct competitor, chain the three inequalities
$Q_h(s,b)\leq U_h(s,b)\leq L_h(s,a)\leq Q_h(s,a)$.
The comparison with $a$ itself is reflexive.
\end{proof}

This is \path{competitor_dominance_sound} in
\path{AVOC/CalibrationBridge.lean}. Excluding self-comparison is essential
for a useful non-singleton certificate. The older sufficient condition compares
every feasible upper score, including the chosen action's own upper score, to
its lower score. That older condition is sound but forces the chosen radius to
be zero. The formal declaration
\path{self_comparison_forces_zero_radius} records this limitation.

\begin{example}[A deliberately conservative abstract quote fixture]
The retained uncertainty example assigns zero setup to both kinds, warm gross
$200$, cold gross $60$, warm success band $[3/5,1]$, and exact cold success $1$.
Start with both statuses untried and use $h=1$, hence two total actions. The
nominal scores are warm $120$, cold $60$, and stop $0$. Here
$\varepsilon=0$, $\delta=4/5$, $M=260$, and $D_1=0$. The local warm radius
is $208$, while the cold and stop radii are zero:
\[
\begin{array}{c|r|r|r}
 \text{action}&\widehat Q_1&\rho_1&L_1\\\hline
 \mathrm{inspectWarm}&120&208&-88\\
 \mathrm{generateCold}&60&0&60\\
 \mathrm{stop}&0&0&0
\end{array}
\]
The lower-Q rule chooses cold although warm's exact compatible Q range is
$[120,200]$. This is a clear example of conservatism in the certified envelope.
It does not show that cold is the exact maximin choice.
\end{example}

The source proves the widths, local radii, and numerical order through
\path{uncertain_warm_local_radius},
\path{uncertain_cold_local_radius}, and
\path{uncertainty_changes_numerical_order}. The displayed first-action
ranking is a direct reconstruction at the stated initial state and horizon.
This fixture is \emph{abstract quote planning}, not an artifact-backed saving
example. Its protected rank-one source has dense cost $36$ per call and the
fixture has one call, so gross values $200$ and $60$ cannot come from bounded
execution savings for that source. The separate $1000$-call example above is
the artifact-backed comparison. These two roles must remain distinct.

\subsection{Whole-policy regret, not only first-action quality}

Let $a_h(s)$ maximize $L_h(s,a)$ and let $a^*$ maximize the true $Q_h$.
Then
\begin{align*}
 V_{h+1}(s)-Q_h(s,a_h(s))
 &\leq U_h(s,a^*)-L_h(s,a_h(s))\\
 &\leq U_h(s,a^*)-L_h(s,a^*)\\
 &\leq 2\Gamma_h.
\end{align*}
The comparison controls the Bellman loss of the actual choice at each state.
The policy is rebuilt at every remaining horizon with the same frozen evidence.

\begin{theorem}[Expected frozen-policy regret]
For a horizon $H$, let $J_H^{\rm lower}(s)$ be the expected value of the complete
adaptive policy induced by robust lower-Q selection. Under any compatible
realization,
\[
 0\leq V_H(s)-J_H^{\rm lower}(s)
 \leq\sum_{h=0}^{H-1}2\Gamma_h.
\]
The comparator is the true optimal legal adaptive plan with the same
conservative reservation guards.
\end{theorem}
\begin{proof}
Write the regret at a node with $h+1$ remaining steps as its local Bellman loss
plus the expected regret of the selected continuation. The local term is at
most $2\Gamma_h$. By induction, every child regret is bounded by the common
sum $\sum_{j<h}2\Gamma_j$. Taking the expectation leaves that constant unchanged
because probabilities sum to one. This yields the cumulative upper bound.
Nonnegativity follows because the selected complete tree is legal and $V_H$
dominates the value of every legal tree.
\end{proof}

This is the argument implemented by \path{plan_regret_le} in
\path{AVOC/GuardedPolicyRegret.lean} and
\path{robust_policy_regret} in \path{AVOC/RobustQuotePlanning.lean}.
Ordinary finite sums simplify the displayed bound to
\[
 \varepsilon H(H+1)+\frac{\delta M}{3}H(H-1)(H+1).
\]
This last closed form is a prose algebraic consequence of the source's sums,
not a newly compiled Lean declaration. The envelope is deliberately broad and
may be vacuous for a particular application. It remains a valid certificate.
The regret is an expected policy-value difference. It does not bound every
realized sequence of successes and failures.

\section{Exact finite-grid calibration}
\label{sec:calibration}

Probability bands need an evidence rule. The implemented rule is an exact
rational finite-grid cell-envelope construction. It should not be relabeled as
an exact inversion of binomial tails or as a Clopper--Pearson interval. Its proof
uses a different and simpler exclusion budget.

Let $K$ be the number of successes in $n$ Bernoulli trials with rational
parameter $p\in[0,1]$. Define
\[
 b_{n,k}(p)=\binom nk p^k(1-p)^{n-k},\qquad 0\leq k\leq n.
\]
Choose a positive integer grid size $m$ and nonnegative rational allowance
$\alpha$. The grid cells have endpoints
\[
 \ell_i=\frac{i}{m},\qquad u_i=\frac{i+1}{m},\qquad 0\leq i<m.
\]
For observed count $k$, the envelope of its point mass over cell $i$ is
\[
 B_{n,k,i}=\binom nk u_i^k(1-\ell_i)^{n-k}.
\]
Retain the cell exactly when
\[
 B_{n,k,i}>\frac{\alpha}{n+1}.
\]
Return the closed interval spanning all retained cells. If there are no retained
cells, return $[0,1]$. The returned interval is therefore always a valid band.
The construction performs exact rational arithmetic and strict comparison.
Boundary cells cover $0$ and $1$ as well as interior parameters.

\begin{theorem}[Finite-sample grid coverage]
If $K$ has the binomial count law $b_{n,k}(p)$, the returned interval $I(K)$ obeys
\[
 \Prob\{p\notin I(K)\}\leq\alpha.
\]
This holds for every rational $p\in[0,1]$, without assuming that $p$ is a grid
point.
\end{theorem}
\begin{proof}
Some grid cell contains $p$. In that cell, $p\leq u_i$ and
$1-p\leq1-\ell_i$, hence $b_{n,k}(p)\leq B_{n,k,i}$. If $p$ is outside the
returned interval, the containing cell could not have been retained: any
retained cell lies inside the hull. Therefore
$B_{n,k,i}\leq\alpha/(n+1)$, and the mass of this excluded count $k$ is at most
that threshold. There are $n+1$ possible counts, so summing their failure masses
gives at most $\alpha$. If no cells were retained, the returned full band cannot
exclude $p$ and the conclusion is immediate.
\end{proof}

The source proves the grid cover, the cell envelope, excluded-count mass, and
final sum as \path{grid_covers}, \path{mass_le_envelope},
\path{excluded_mass_le}, and \path{missMass_le} in
\path{AVOC/FiniteCalibration.lean}. This proves finite-sample validity of the
algorithm itself within the exact rational binomial model. It does not infer the
sampling model from a list of reported outcomes.

\begin{example}[Ten successes are evidence, not certainty]
For $n=10$, $m=10$, and $\alpha=1/10$, ten successes retain cells
$\{6,7,8,9\}$ and return $[3/5,1]$. Ten failures retain
$\{0,1,2,3\}$ and return $[0,2/5]$. These endpoint calculations are explicit
Lean theorems. By contrast, the empirical singleton $\{K/n\}$ has no such
coverage property. At $n=1$ and $p=1/2$, both possible singletons exclude the
true parameter, so its miss probability is $1$.
\end{example}

The relevant declarations include \path{ten_successes_lower},
\path{ten_successes_upper}, \path{ten_failures_lower},
\path{ten_failures_upper}, and
\path{empirical_singleton_counterexample}. Nothing here says the true
parameter has posterior probability $9/10$ of lying inside an observed band.
The guarantee concerns repeated-data failure probability under the specified
law.

\subsection{Predetermined schedules and data-dependent selection}

A schedule is a finite set $J$ of checkpoints. Each checkpoint fixes an
identity stratum, sample size $n_j$, grid $m_j>0$, true parameter $p_j$, and
nonnegative error allocation $\alpha_j$. Its count statistic has the stated
binomial marginal. Counts can be dependent, including different prefixes of
one source and checkpoints that share other underlying randomness.

\begin{theorem}[Scheduled and selected failure bounds]
Let $B_j$ be the event that checkpoint $j$'s interval excludes $p_j$. Then
\[
 \Prob\Bigl(\bigcup_{j\in J}B_j\Bigr)
 \leq\sum_{j\in J}\alpha_j.
\]
For any data-dependent selector $j(\omega)\in J$,
\[
 \Prob\{\omega\in B_{j(\omega)}\}\leq\sum_{j\in J}\alpha_j.
\]
\end{theorem}
\begin{proof}
For each outcome $\omega$, the indicator of a union is at most the sum of the
individual indicators. Multiply by the nonnegative rational weight of
$\omega$ and sum over the finite probability space. Interchange the two finite
sums to obtain the union bound. Each checkpoint has failure mass at most its
allocated $\alpha_j$ by finite-grid coverage. Finally, selected failure implies
failure of at least one scheduled checkpoint, so it is a subset of the union.
This last inclusion holds pointwise even if selection sees all outcomes.
\end{proof}

This reconstructs \path{finite_union_le},
\path{scheduled_binomial_failure_le}, and
\path{selected_binomial_failure_le} in
\path{AVOC/CalibrationSchedule.lean}. Neither checkpoint independence nor
an independent selector is required. Predetermination is required: the theorem
does not authorize creating arbitrary new checkpoints after viewing the data.
It is finite scheduled selection, rather than an unrestricted anytime
confidence sequence.

\begin{example}[Why the schedule's total allowance is necessary]
Take a fair Boolean outcome $\omega$ and two checkpoint bad events
$B_0=\{\omega=0\}$ and $B_1=\{\omega=1\}$. Each has probability $1/2$.
Selecting $j(\omega)=\omega$ makes selected failure certain. The total allowance
$1/2+1/2=1$ remains valid, while either single checkpoint's allowance fails.
This is the exact obstruction proved by
\path{fairBool_adaptive_exceeds_single_budget}.
\end{example}

\section{From source laws to actual terminal samples}
\label{sec:recorder}

The count-law obligation is substantial. Fair marginals do not imply a binomial
sum. Two identical fair bits each have parameter $1/2$, but their two-bit count
is either zero or two and never one. The source contains this obstruction as
\path{correlatedPair_not_product}.

The primitive model uses a finite rational probability space and, for each
identity $k$, a finite latent Boolean vector of length $N_k$. Its success
parameter $p_k$ is fixed. The fresh-outcome contract states that after agreement
with any specified pattern on a previously revealed set of positions, a fresh
position's joint success mass is $p_k$ times that history mass. In notation,
\[
 \Prob\{H\text{ matches and }X_i=1\}=p_k\Prob\{H\text{ matches}\},
 \qquad i\text{ fresh}.
\]
The multiplication-only form is meaningful for histories of mass zero and does
not divide by them. Induction determines the mass of every full pattern, then
summing patterns by success count gives its binomial law. Truncation yields the
same law for each fixed prefix.

The formal steps are \path{conditional_implies_product},
\path{count_fiber_eq_binomial}, and
\path{prefix_count_fiber_eq_binomial} in
\path{AVOC/ObservationStream.lean}. This derives the count law from a
primitive source law. The primitive source law itself remains an explicit
model assumption.

\subsection{Pending intents and immutable terminal events}

The new lifecycle recorder separates a pending-intent list from a terminal log.
A single global natural-number operation-ID namespace covers both lists and
all identities. Pending entries store an ID and identity. Terminal entries
store an ID, identity, and success, rejection, or fault. An unfinished outcome
cannot be inserted as a terminal event.

\begin{center}
\begin{tabularx}{\linewidth}{@{}lX@{}}
\toprule
Command situation & Defined result\\\midrule
Fresh begin & Append one pending ID and identity.\\
Exact pending begin replay & Return the same state.\\
Begin replay after matching completion & Return the same state without reopening.\\
Matching fresh finish & Remove the pending ID and append one terminal event.\\
Exact terminal replay & Return the same state.\\
Changed identity or terminal outcome & Reject the command.\\
Finish without matching intent & Reject the command.\\
Unfinished terminal outcome & Reject the command.\\
\bottomrule
\end{tabularx}
\end{center}

These cases resolve a concrete inherited obstruction. If an immutable log stores
an unfinished event under an ID, later adding success under that ID changes the
event and must be rejected. A pending registry therefore cannot simply be an
unfinished terminal record. The declaration
\path{inherited_unfinished_terminal_conflict} proves the mismatch, rather
than relying on a narrative explanation.

\begin{theorem}[Pure recorder lifecycle invariant]
Every successful run from the empty state has unique pending IDs, unique
terminal IDs, disjoint pending and terminal IDs, and no unfinished terminal
entries. An accepted terminal event occurs once and remains once throughout
every later successful continuation. Terminal history grows by suffix append.
\end{theorem}
\begin{proof}
For a begin step, the transition first checks both uses of the ID. It either
returns the unchanged state for a matching replay, rejects conflicting reuse,
or appends a fresh pending entry. For a finish step, exact replay is unchanged;
otherwise successful completion requires an unused terminal ID and a matching
pending entry. Removing that pending ID and appending its terminal event
preserves uniqueness and disjointness. Unfinished outcomes are rejected before
these cases. Induction over a successful command list preserves the invariant.
Every successful step leaves the log unchanged or appends one event, which
proves retention under continuation.
\end{proof}

The source declarations are \path{run_wellFormed},
\path{runFrom_log_prefix}, and
\path{accepted_terminal_retained_once} in
\path{AVOC/ObservationRecorder.lean}. A success adds true once to its
identity's samples, a rejection adds false once, and a fault adds no sample.
False outcomes count toward sample readiness just as true outcomes do.

\subsection{A source-generative recorder and whole-run refinement}

The raw recorder accepts externally supplied terminal labels. The generative
recorder instead has begin, report, and fault commands. A fresh report constructs
its true or false label from the next latent Boolean of its identity. A replay
checks the terminal log first and reuses the stored outcome, consuming no new
source bit. A fresh report with an exhausted source rejects. Faults and begins
consume no source bit. Reporting after a fault conflicts with the immutable
terminal event.

Every accepted generative step emits an ordinary recorder command and an
effective reveal marker. Whole-run induction proves that the emitted commands
reproduce the actual raw recorder state and that the reveal state equals the
inherited reveal machine on its effective trace. In particular, for every
identity,
\begin{equation}
 \mathrm{samples}_k(\mathrm{terminalLog})
 =\mathrm{source}_k.\mathrm{take}
     \bigl(|\mathrm{samples}_k(\mathrm{terminalLog})|\bigr).
 \label{eq:recorder-prefix}
\end{equation}

This is \path{run_refines} and \path{run_samples_eq_take} in
\path{AVOC/RecorderReveal.lean}. The source-prefix correspondence is a
derived invariant of the defined machine, not an additional equality asserted
for each observed raw event. It still does not establish that an arbitrary
external success label came from this latent source. A parser can check the
label's syntax without checking its physical or statistical origin.

\begin{example}[Replay and faults do not create new evidence]
Let an identity's source be $[\mathrm{true},\mathrm{false}]$. Begin ID $7$, report
ID $7$, then report ID $7$ again. The sample list is $[\mathrm{true}]$ and the
remaining source begins with false. Begin ID $8$ and finish it as a fault. The
sample list stays unchanged. Begin and report ID $9$. It appends false, yielding
$[\mathrm{true},\mathrm{false}]$. A further fresh report rejects after source
exhaustion, while a replay of ID $7$ remains possible. The replay supplies no
extra success and the fault supplies no concealed binary observation.
\end{example}

This is an illustrative execution of the defined transitions. The retained
formal examples \path{replay_fault_exhaustion_example} and
\path{exhausted_fresh_report_rejected} provide closely related checked
instances. Reusing a different ID for the same physical operation is a separate
identity problem. Global ID uniqueness alone cannot rule it out.

\section{Unconditional calibrated policy issuance}
\label{sec:issuance}

The recorder-to-policy endpoint reads counts from the actual terminal log.
Rejected whole runs supply no counts. Successful runs with \emph{any} pending
operation also supply no counts, even when the pending identity is unrelated to
the selected checkpoint. A fully resolved run still needs enough success or
rejection reports for both selected checkpoints. A checkpoint always reads the
first predetermined $n_j$ reports, so later reports cannot change that count.

When both selected counts are available, the endpoint constructs their exact
learned probability bands and installs them for a fresh frozen policy episode.
Setup and gross memberships are supplied deterministic premises. The warm and
cold selectors may depend on all outcomes but must select checkpoints with the
specified warm and cold true parameters. Changing a probability band does not
authorize changing an outstanding operation's evidence.

\begin{theorem}[Recorder-to-policy bad-issuance bound]
Assume finite rational latent sources satisfy the primitive conditional laws;
the finite schedule, grids, and nonnegative allowances are predetermined;
scheduled prefix lengths fit within their sources; selected strata match the
specified true probabilities; and deterministic setup/gross memberships hold.
Then
\begin{equation}
 \Prob\{\text{a plan is issued and lacks its episode certificate}\}
 \leq \sum_{j\in J}\alpha_j.
 \label{eq:issuance}
\end{equation}
The episode certificate supplies a compatible realization with the stated exact
parameters, value and Q enclosures, and expected whole-policy regret for every
finite horizon under frozen quotes.
\end{theorem}
\begin{proof}
The primitive conditional law yields binomial prefix count laws at all scheduled
checkpoints. The generative run refinement gives \eqref{eq:recorder-prefix}.
Whenever gated counts are available, taking the fixed first $n_j$ reports
therefore gives the corresponding latent prefix count. On the event where all
scheduled bands cover their true parameters, either the gates suppress issuance
or both selected learned bands contain the correct parameters. Combine those
memberships with the supplied deterministic memberships to construct a
compatible realization. The deterministic frozen-episode theorem then supplies
the policy certificate. Thus bad issuance is a subset of the scheduled band
failure union, whose probability is bounded by the total allowance.
\end{proof}

The final declarations are \path{recorder_episode_coverage} and
\path{recorder_issued_failure_le} in
\path{AVOC/RecorderPolicy.lean}. Success of execution, absence of pending
intents, availability of enough samples, log-prefix equality, and the binomial
count law are not separately supplied premises to that final theorem. They are
derived or checked by its machine and gates. The primitive source laws and
deterministic memberships are still supplied.

The probability in \eqref{eq:issuance} ranges over the original outcome space,
including rejected, unresolved, and insufficient-data runs. There is no claim
that a plan will be issued, that any operation finishes, or that collection is
cheap. If the issuance probability is $q>0$, the generic conditional implication
is only
\[
 \Prob\{\text{uncertified}\mid\text{issued}\}
 \leq \min\left(1,\frac{\sum_j\alpha_j}{q}\right).
\]
The same nominal allowance need not hold after conditioning on issuance.

\subsection{Completion conditioning can select a biased population}

\begin{example}[Predictable availability still biases completed runs]
Let $(X_1,X_2)$ be independent fair bits. Reveal $X_1$, and reveal $X_2$ by the
deadline only if $X_1=1$. Availability of the second report depends on the past,
not on $X_2$. Completion probability is $1/2$. Conditional on completion,
$X_1=1$ with probability one, and the two-bit count has probability zero of
being zero. A genuine two-trial binomial count with parameter $1/2$ has zero
count probability $1/4$.
\end{example}

This is kernel-checked in \path{AVOC/ObservationCensoring.lean}, including
\path{deadline_second_retention_predictable},
\path{deadline_conditioned_first_success}, and
\path{deadline_conditioned_count_not_binomial}. It shows why one cannot
replace the original law by the law conditioned on enough reports arriving.

A direct interval illustration uses ten independent fair bits and issues an
interval only when all ten are true. Its issuance probability is $2^{-10}$.
The implemented $n=m=10$, $\alpha=1/10$ interval is then $[3/5,1]$, which excludes
$1/2$. Conditional on issuance, failure is certain. Unconditional bad issuance
is only $2^{-10}$ and satisfies the allowance. This finite calculation is an
illustrative consequence of the checked interval endpoints and the elementary
product law, not a new Lean declaration.

\subsection{Complete logs do not fix outcome-dependent completion order}

Conditioning and ordering are different obstructions. Suppose two independent
fair attempts start together. Success takes one unit of time and rejection
takes two; resolve ties by the initiation ID. Record every final outcome and
wait until all intents have resolved. The completion order is
\[
\begin{array}{c|c|c}
\text{initiation outcomes}&\text{completion-order outcomes}&\text{mass}\\\hline
(0,0)&(0,0)&1/4\\
(0,1)&(1,0)&1/4\\
(1,0)&(1,0)&1/4\\
(1,1)&(1,1)&1/4
\end{array}
\]
The first completed bit succeeds with probability $3/4$. No report was missing,
no terminal event was forged, and no operation remains pending. Nevertheless,
the first-completed prefix is not a fair source prefix.

For the implemented interval, the stronger retained finite experiment starts
twenty independent fair attempts and lets all successes complete before all
rejections. If there are at least ten successes, the first ten completed reports
are all true. Their interval is $[3/5,1]$ and excludes $1/2$. That event has mass
\[
 \frac{1}{2^{20}}\sum_{k=10}^{20}\binom{20}{k}
 =\frac{308333}{524288}>\frac{1}{10}.
\]
Thus actual interval failure exceeds its nominal allowance despite complete,
fully resolved lifecycle records. The exact Python evidence is preserved under
\path{results/v0.41.0/review_completion_order.py} and its JSON receipt. It is
a finite computational obstruction with an explicit rational derivation here,
not a newly kernel-checked theorem.

The generative theorem assigns the next latent bit in report order and assumes
that source's conditional law. Independent outcomes in initiation order do not
automatically satisfy that assumption after outcome-dependent reordering. A
physical deployment must justify its report-order contract, for example through
an initiation-indexed sampling policy or a proved alternative coupling. The
current release leaves that physical coupling open.

\section{Checked recovery and durable accounting}
\label{sec:recovery}

Recovery must distinguish a recorded intention from a completed paid result.
The robust quote recovery model binds a commit to its sequence, complete
pre-state, selected action, quote evidence, outcome tag, and post-state. Fresh
acknowledgment checks those bindings and checks that the supplied typed outcome
realizes the action through the runtime's \path{confirm} transition. It then
requires the computed successor's store view to equal the recorded durable
post-state. An unfinished intent is rejected rather than interpreted as a
negative observation or permission to repeat paid work.

\begin{proposition}[Fresh acknowledgment and replay]
A successful fresh acknowledgment advances the sequence by one and yields
the successor certified by \path{confirm}. It preserves the runtime safety
invariant. Repeating that accepted commit at the resulting matching checkpoint
returns the same checkpoint.
\end{proposition}
\begin{proof}
Each failed identity, sequence, pre-state, outcome, evidence, or post-state
comparison returns an error. A successful fresh branch therefore exposes all
those equalities and a successful \path{confirm}. The successor invariant
follows from \path{confirm_safe}. At the resulting checkpoint, the recorded
last commit, next sequence, and actual store view match the duplicate branch,
which returns the state unchanged.
\end{proof}

These are \path{acknowledge_fresh_contract},
\path{acknowledge_safe}, and \path{acknowledge_idempotent} in
\path{AVOC/RobustQuoteRecovery.lean}. Duplicate acknowledgment recognizes
the already accepted commit and checkpoint. It does not freshly validate every
field of an independently supplied outcome argument on that replay branch.
This is appropriate for state replay but must not be described as new checking
of a second physical result.

The \path{prepare} constructor selects the robust action and binds current
evidence into its intent. Acknowledgment alone does not prove that an arbitrary
supplied commit was originally produced by \path{prepare}. If the conclusion
needs provenance of the planner's choice, that provenance must be supplied or
established by the enclosing protocol. Fresh acknowledgment proves the exact
checked transition contract, not an unstated origin claim.

Probability bands are frozen within the episode. Replacing them for a fresh
episode preserves the source and protected identity fields but changes the
approved evidence. An old intent cannot be acknowledged under the changed
evidence. The explicit declaration
\path{new_bands_obstruct_old_acknowledgement} in
\path{AVOC/CalibrationBridge.lean} prevents calibration from silently
rewriting an active durable operation.

\subsection{Prepaid evaluation and upgrade journals}

The separate durable-upgrade machine has a journal containing the current
upgrade state, one optional pending ticket, and closed tickets. A ticket binds
a campaign-local key, its complete parent state, and a proposal. Beginning a
funded fresh ticket charges its protected evaluation tariff immediately.
Repeating the same ticket is unchanged. Finishing it charges no second
evaluation tariff, and switching occurs only according to the checked upgrade
decision. Abandoning clears the pending ticket while retaining its already
charged evaluation expenditure.

The upgrade authority supplies checker soundness and sound lower and upper
utility bounds for one fixed workload, objective, and unit convention. These
are explicit obligations in \path{AVOC/VerifiedUpgrade.lean}. Receiving a JSON
utility bound or an approval flag does not prove them.

In particular, a failed or abandoned evaluation is not a free trial. The
complete parent state includes spending, so reusing a ticket based on an old
active version but ignoring changed expenditure is rejected. The formal
declarations \path{begin_keeps_active},
\path{finish_evaluation_unchanged}, and
\path{abandon_keeps_current} make these distinctions explicit.

For a successful paid trace with total evaluation debit $e$ and switching debit
$u$, the formal accounting identities are
\begin{align}
 \mathrm{evaluationSpent}_{\rm final}
 &=\mathrm{evaluationSpent}_{\rm initial}+e,\\
 \mathrm{switchSpent}_{\rm final}
 &=\mathrm{switchSpent}_{\rm initial}+u.
\end{align}
From a safe initial journal, their sum stays within protected credits.
Successful replay also cannot decrease the protected utility of the active
version. These are \path{paid_trace_exact_accounting},
\path{paid_trace_no_overspend}, and
\path{replay_utility_nondecreasing} in
\path{AVOC/DurableUpgradeLaws.lean}.

Restoration replays supplied commands from a specified initial version and
compares the result with the actual checkpoint. A successful restoration equals
that checkpoint. Provided the specified initial version satisfies the
authority's contract, restoration also inherits the safety certificate. This checks consistency
between two supplied records. It cannot authenticate their physical origin or
detect coordinated replacement of both the journal and checkpoint.

\subsection{What recovery establishes about external execution}

The lifecycle recorder is a pure finite snapshot reducer. Its once-only
terminal history is not exactly-once external execution. It neither invokes
an operation nor persists a crash-safe transaction. The recovery models give
checked state-transition and replay contracts, while the retained runtime and
storage adapters have their stated bounded evidence and external trust
assumptions. Filesystem durability, endpoint authentication, physical identity,
and reliable observation delivery are not consequences of a pure transition
function.

A rejected pure batch supplies no final export. That is an all-or-nothing
reduction result, not rollback of physical work already performed before
rejection. A fault closes an intent without supplying a binary sample. It does
not establish that retrying the corresponding paid operation is safe or free.
Conservative reserved spending remains a protected model quantity rather than
an automatic identity with actual billed expenditure.

The completed chain is therefore precise: under the declared finite source
model, the lifecycle machinery derives trustworthy prefix correspondence;
fixed scheduled calibration bounds bad issuance; on the simultaneous-coverage
event, every issued frozen episode has the proved value and policy guarantees;
outside that event, uncertified issuance has mass at most the scheduled total
allowance; and the checked recovery machines preserve their specified state
and accounting contracts. Physical source coupling,
universal host/storage refinement, multi-episode accumulation, and measured
productivity remain separate obligations.
